% ANCHOR: after_label def:c8-twocomp
\begin{remark}[Physics notation: ten-dimensional and $SO(6)$ gamma matrices]
With $\zeta_A=(\xi_\alpha{}^i,\bar\xi_{\dot\alpha i})\in\Sigma=\mathbf{16}$ and $\chi^A=(\eta^\alpha{}_i,\bar\eta^{\dot\alpha i})\in\Sigma^*=\mathbf{16}'$ of $\mathrm{Spin}(10)$ (the key after Section 1.2), the gamma matrices $\Gamma_m=\rho(e_m)$ are the formulas of the definition at $v=e_m$, and the six transverse ones $\Gamma_I=\rho(e_I)$ are, with $\bar\Sigma^I=(\Sigma^I)^\dagger$ and $w^{ij}=-w_I\bar\Sigma^{I\,ij}$,
\[
\begin{aligned}&(\Gamma_I\zeta)^\alpha{}_j=\varepsilon^{\alpha\beta}\Sigma^I_{ji}\xi_\beta{}^i,\qquad (\Gamma_I\zeta)^{\dot\alpha j}=-\varepsilon^{\dot\alpha\dot\beta}\bar\Sigma^{I\,ji}\bar\xi_{\dot\beta i},\\ &(\Gamma_I\chi)_\alpha{}^k=\bar\Sigma^{I\,kj}\varepsilon_{\alpha\beta}\eta^\beta{}_j,\qquad (\Gamma_I\chi)_{\dot\beta k}=-\Sigma^I_{kj}\varepsilon_{\dot\beta\dot\alpha}\bar\eta^{\dot\alpha j}.\end{aligned}
\]
The 't Hooft blocks are the chiral blocks of the $SO(6)_R$ gamma matrices: on $W\oplus W^*$,
\[
\gamma^I=\begin{pmatrix}0&\bar\Sigma^I\\ \Sigma^I&0\end{pmatrix},\qquad \gamma^I\gamma^J+\gamma^J\gamma^I=2\delta^{IJ},\qquad (\gamma^I)^\dagger=\gamma^I,
\]
so $W=\mathbf4$ and $W^*=\bar{\mathbf4}$ are the two Weyl spinors of $\mathrm{Spin}(6)\cong SU(4)_R$, and $W_{\mathbb R}^\perp$ is its vector $\mathbf6$. The generators $Q^i_\alpha$, $\bar Q_{\dot\alpha i}$, $S^\alpha{}_i$, $\bar S^{\dot\alpha i}$ span $(\mathbf2,\mathbf1,\mathbf4)$, $(\mathbf1,\mathbf2,\bar{\mathbf4})$, $(\mathbf2,\mathbf1,\bar{\mathbf4})$, $(\mathbf1,\mathbf2,\mathbf4)$ of $SU(2)_L\times SU(2)_R\times SU(4)_R$.
\end{remark}

% ANCHOR: after_label prop:c8-twocomp
\begin{remark}[Physics notation: Feynman gauge with two cutoffs]
The propagator is that of Feynman gauge. As in Proposition \ref{prop:c3-P}, $\hbar$ times the inverse pairing is $\tfrac12\delta^{ab}$ in the basis $-igT^a$ and $\delta^{ab}$ in the basis $-\tfrac{ig}{\sqrt2}T^a$, so the edges $\hbar P(\varepsilon,T)$ of the flow are the canonical Euclidean propagators
\[
\hbar\langle A^a_mA^b_n\rangle=\delta^{ab}\delta_{mn},\qquad \hbar\langle A^a_IA^b_J\rangle=\delta^{ab}\delta_{IJ},\qquad \hbar\langle\lambda^a_\alpha{}^i\,\bar\lambda^b_{\dot\beta j}\rangle=-\delta^{ab}\delta^i_j\,p_m\sigma^m_{\alpha\dot\beta},
\]
each times $\big(e^{-\varepsilon|p|^2}-e^{-T|p|^2}\big)/|p|^2$, with the proper-time cutoffs $\varepsilon$ (ultraviolet) and $T$ (infrared). In $SO(6)$ form, $\langle A_{ij}A_{kl}\rangle=4(-,-)^{-1}\varepsilon_{ijkl}$ is $\langle A_IA_J\rangle=2(-,-)^{-1}\delta_{IJ}$, by $A_I=\tfrac14\bar\Sigma^{I\,ji}A_{ij}$.
\end{remark}

% ANCHOR: end_of_section 8.1
\begin{remark}[Physics notation: odd brackets on unit elements]
On the unit elements of Definition \ref{def:c8-twocomp}, the odd--odd brackets of Chapter 2 (graded brackets, $P_m=\partial_m$) are
\[
[Q^i_\alpha,\bar Q_{\dot\alpha j}]=i\,\delta^i_j\,\bar\sigma^{m\dot\alpha\alpha}P_m,\qquad [Q^i_\alpha,Q^j_\beta]=0,\qquad [\bar Q_{\dot\alpha i},\bar Q_{\dot\beta j}]=0,
\]
the values of $[\xi_1,\xi_2]=i\big(\bar\xi_{1\dot\alpha i}\bar\sigma^{m\dot\alpha\alpha}\xi_{2\alpha}{}^i+\bar\xi_{2\dot\alpha i}\bar\sigma^{m\dot\alpha\alpha}\xi_{1\alpha}{}^i\big)P_m$, the Euclidean form of the anticommutator of $Q$ and $\bar Q$ into a translation ($\bar\sigma^mP_m=\bar\sigma^\mu P_\mu$ after the Wick map). The dilatation part of $[Q^i_\alpha,S^\beta{}_j]$ is $\tfrac12\delta_\alpha{}^\beta\delta^i{}_j\Delta$ and that of $[\bar Q_{\dot\alpha i},\bar S^{\dot\beta j}]$ is $-\tfrac12\delta_{\dot\alpha}{}^{\dot\beta}\delta_i{}^j\Delta$, from the term $\tfrac12(\xi,\eta)\Delta$ of $[\xi,\rho(x)\eta]$; the brackets of $Q$ with $\bar S$ and of $\bar Q$ with $S$ have no dilatation part.
\end{remark}

% ANCHOR: end_of_section 8.2
\begin{remark}[Physics notation: infrared endpoint and Wilsonian flow]
$T$ is an infrared proper-time cutoff, of length $T^{1/2}$, and $\varepsilon$ the ultraviolet one: the scalar factor of $P(\varepsilon,T)$ is the massless propagator $\frac1{4\pi^2|x-y|^2}$ with both cutoffs, $\frac{e^{-|x-y|^2/4T}-e^{-|x-y|^2/4\varepsilon}}{4\pi^2|x-y|^2}$. $W(P(\varepsilon,T),J)$ is the Wilsonian effective interaction at scale $T$, the connected Feynman graphs with vertices from $J/\hbar$ and internal lines $\hbar P(\varepsilon,T)$, in which the modes of proper time between $\varepsilon$ and $T$ are integrated out. Lemma \ref{lem:c8-transfer} is the semigroup property of this exact renormalisation-group flow and its compatibility with the BV differential.
\end{remark}

% ANCHOR: after_label def:c8-quant
\begin{remark}[Physics notation: the equivariant quantum master equation]
$I_{\mathfrak{psl}(4|4)}[T]$ is the Wilsonian effective action at scale $T$ with the constant superconformal ghosts. \eqref{eq:c8-qme} is the regularised Slavnov--Taylor identity: its arity-zero part is that of BRST invariance, its arity-one part along $x$ the Ward identity of the generator $x$, its arity-two part the closure of the algebra. $F_{\mathfrak{psl}(4|4)}$ is a finite symmetry-restoring counterterm in the sense of algebraic renormalisation, and the functions $\varepsilon^{-m}\log^k\varepsilon$ are the power and logarithmic divergences of the proper-time cutoff.
\end{remark}

% ANCHOR: after_label thm:c8-obstruction
\begin{remark}[Physics notation: anomaly and renormalisation ambiguities]
$\mathrm{Ob}$ is the one-loop anomaly of the regularised Slavnov--Taylor and superconformal Ward identities; (i) is its Wess--Zumino consistency condition, (ii) states that it is removed by a finite counterterm exactly when it is $s$-exact, and (iii) identifies the finite renormalisation ambiguities with $H^0$ of $s$ on local cochains, the field redefinitions being trivial. In (iv) the coupling direction is a finite renormalisation of the coupling: through the pairing $(X,Y)=-\tfrac{4\hbar}{g^2}\operatorname{tr}(XY)$ the action $S$ is proportional to $g^{-2}$ at fixed $\mathfrak g$-valued fields, and $\mu S$ is $S$ with $g^{-2}\mapsto\mu g^{-2}$. $G_B$ is the relative normalisation of $Q$ and $\bar Q$ (Theorem \ref{thm:c8-invariants} (ii)).
\end{remark}

% ANCHOR: after_label def:c8-B
\begin{remark}[Physics notation: the hypercharge]
$B$ is the $\mathfrak u(1)$ outer automorphism of $\mathfrak{psu}(2,2|4)$, the hypercharge $U(1)_Y$: in $\mathfrak{gl}(4|4)$, with real form $\mathfrak u(2,2|4)$, it is the commutator with $\operatorname{diag}(0_4,1_4)$. It vanishes on $P$, $M$, $\bar M$, $\Delta$, $K$, $R$ and counts upper minus lower $SU(4)_R$ indices on the supercharges,
\[
B(Q^i_\alpha)=Q^i_\alpha,\qquad B(\bar Q_{\dot\alpha i})=-\bar Q_{\dot\alpha i},\qquad B(S^\alpha{}_i)=-S^\alpha{}_i,\qquad B(\bar S^{\dot\alpha i})=\bar S^{\dot\alpha i},
\]
so $Q^i_\alpha$ and $\bar S^{\dot\alpha i}$ (in the $\mathbf4$) have charge $+1$ and $\bar Q_{\dot\alpha i}$ and $S^\alpha{}_i$ (in the $\bar{\mathbf4}$) charge $-1$. In components $G_B^{(1)}(x)=S^{(1)}(Bx)$ is $S^{(1)}$ with $\xi_\alpha{}^i$, $\bar\eta^{\dot\alpha i}$ kept and $\bar\xi_{\dot\alpha i}$, $\eta^\alpha{}_i$ reversed in sign.
\end{remark}

% ANCHOR: end_of_section 8.3
\begin{remark}[Physics notation: the two finite ambiguities]
$e^{tB}$ multiplies $Q^i_\alpha$, $\bar S^{\dot\alpha i}$ by $e^t$ and $\bar Q_{\dot\alpha i}$, $S^\alpha{}_i$ by $e^{-t}$ and fixes the even part, an automorphism of the superconformal algebra; $G_B$ is the corresponding change of the couplings to the constant ghosts, with $G_B^{(0)}=0$, so the action $S$ itself is unchanged. Proposition \ref{prop:c8-B} states that this ambiguity, a relative rescaling of $Q$ and $\bar Q$, is independent of the finite renormalisation of the coupling.
\end{remark}

% ANCHOR: after_label def:c8-components
\begin{remark}[Physics notation: the one-loop generators]
$Q_1(O)^\alpha{}_i$, $\bar Q_1(O)^{\dot\alpha i}$, $S_1(O)_\alpha{}^i$, $\bar S_1(O)_{\dot\alpha i}$ are the one-loop corrections to the transformations of $O$ under $Q^i_\alpha$, $\bar Q_{\dot\alpha i}$, $S^\alpha{}_i$, $\bar S^{\dot\alpha i}$; $P_1$, $M_1$, $R_1$, $K_1$, $D_1$ are those under $P_b$, $\mathfrak{so}(4)$, $\mathfrak{su}(4)_R$, $K_b$, $\Delta$; $\hbar L_1(O)^{(0)}$ is the one-loop correction to the BRST variation of $O$. For the dilatation $D=-\tfrac12\Delta$ of physics, acting through $-\Lambda^{(1)}$ (Corollary \ref{cor:c8-homotopy}), the one-loop dilatation operator is $\tfrac12D_1$: $\tfrac12D_1(O)=-24$ units times $O$, i.e. $\frac{3g^2N}{4\pi^2}O$, for the Konishi operator, and $0$ on the $\mathbf{20'}$ (Remark \ref{rem:c6-konishi}).
\end{remark}

% ANCHOR: end_of_section 8.4
\begin{remark}[Physics notation: renormalised operators and their transformations]
$M(\varepsilon)$ is the bare operator with its counterterms, $\psi$ the finite part that fixes the renormalisation scheme of the composite operators ($\psi=0$ is minimal subtraction of the singular terms $\varepsilon^{-m}\log^k\varepsilon$), and $M[T]$ the renormalised operator in the Wilsonian theory at scale $T$. $d_T$ is the quantum BRST differential at scale $T$, and $L(M)$ the local composite operator representing the quantum transformation of $M$, unique by Theorem \ref{thm:c8-L} (i). The cohomology of $\Lambda^{(0)}$ on local functionals is the space of physical local operators (BRST-closed modulo BRST-exact), and by Corollary \ref{cor:c8-homotopy} $x\mapsto-\Lambda^{(1)}(x)$ is the action of the superconformal algebra on it, closed up to BRST-exact terms; in this normalisation $D=-\tfrac12\Delta$ acts on an operator of degree $0$ at the origin as multiplication by its scaling dimension $w_O+\gamma_O$.
\end{remark}

% ANCHOR: end_of_section 8.5
\begin{remark}[Physics notation: scheme dependence]
(a) and (b) are field redefinitions and finite renormalisations $O\mapsto O+\hbar\psi'(O)$ of the composite operators, changes of scheme that conjugate all one-loop generators. (c) with the coupling direction is a finite renormalisation of $g$, which rescales the one-loop BRST operator, and with $G_B$ the hypercharge rescaling of the supercharges, $Q^i_\alpha$ and $\bar S^{\dot\alpha i}$ by $1-a\hbar$, $\bar Q_{\dot\alpha i}$ and $S^\alpha{}_i$ by $1+a\hbar$. Theorem \ref{thm:c8-invariants} (iii) is the scheme independence of the one-loop anomalous dimensions, the eigenvalues of $\tfrac12D_1$ on operators of fixed canonical dimension.
\end{remark}

% ANCHOR: after_label lem:c8-ev
\begin{remark}[Physics notation: the twisting supercharge]
$[\xi,\xi]=0$ for $\xi\in S_+\otimes W$ is $[Q^i_\alpha,Q^j_\beta]=0$. For $\xi=u\otimes z$ the ten-dimensional spinor $\zeta_A=(u_\alpha z^i,0)$ is a pure spinor, $\Gamma^M(\xi,\xi)=0$ for all ten $M=(m,I)$. For $z^i=\delta^i{}_4$ and $u_\alpha$ a weight vector of $\mathfrak{su}(2)_L$, $\xi$ is the single supercharge written $Q^4_-$ in the physics literature on the holomorphic twist; every nonzero $u\otimes z$ is of this form after an $SU(2)_L\times SU(4)_R$ rotation and a rescaling. $r_\xi$ sets the constant ghost of $\xi$ equal to $q$, then to $1$: $S^\xi$ is the BV action of the theory twisted by $Q^4_-$, its differential $\{S^\xi,-\}$ is BRST plus $Q^4_-$, and
\[
-\tfrac12S^{(2)}(\xi,\xi)=\tfrac14\int\big(u_\alpha z^i\lambda^{*\alpha}{}_i,\ u_\beta z^j\lambda^{*\beta}{}_j\big)
\]
is the term quadratic in the gaugino antifields that is required because $\{S^{(1)}(\xi),S^{(1)}(\xi)\}=\{S,S^{(2)}(\xi,\xi)\}$ is a nonzero multiple of the gaugino equation of motion: $Q^4_-$ squares to zero on shell.
\end{remark}

% ANCHOR: end_of_section 8.6
\begin{remark}[Physics notation: the twisted one-loop differential]
$Q^\xi_0=\{S^\xi,-\}$ is the classical twisted differential, BRST plus $Q^4_-$; its cohomology on local operators is the classical $Q^4_-$-cohomology, the $\tfrac1{16}$-BPS sector of the physics literature. $\hbar Q^\xi_1$ is its one-loop correction, the one-loop BRST correction $\hbar L_1^{(0)}$ plus the one-loop $Q^4_-$-component $Q_1$ plus the arity-two terms, and $[Q^\xi_1]$ is the scheme-independent one-loop correction of $Q^4_-$ on this cohomology, the object compared with the one-loop differential of [BK].
\end{remark}

% ANCHOR: end_of_section 8.7
\begin{remark}[Physics notation: infrared behaviour]
The endpoint $T$ is an infrared regulator. The insertions of $S^{(1)}$ along $Q$, $S$, $\Delta$, $P$, $M$, $\bar M$, $R$ are infrared finite, and the special conformal $K_b$, whose Killing field $x^2b^m-2(b\cdot x)x^m$ grows quadratically, produces an infrared logarithm, $\int_{|y|^2<T}|y|^{-4}d^4y$ growing as $\pi^2\log T$. Theorem \ref{thm:c8-K} states that the local operator $K_1(O)$ is nevertheless infrared safe, the $\log T$ residing in the Wilsonian flow of $O[T]$; in radial quantisation $K_b$ is the conjugate of $P_b$ (Corollary \ref{cor:c8-gram}).
\end{remark}

% ANCHOR: after_label lem:c8-cc
\begin{remark}[Physics notation: Euclidean conjugation of spinors]
In Euclidean signature complex conjugation keeps the four-dimensional chirality and exchanges $\mathbf4$ and $\bar{\mathbf4}$: $\overline{\xi_\alpha{}^i}$, placed as $\eta^\alpha{}_i$, lies in $(\mathbf2,\mathbf1,\bar{\mathbf4})$, and $-\overline{\bar\xi_{\dot\alpha i}}$, placed as $\bar\eta^{\dot\alpha i}$, in $(\mathbf1,\mathbf2,\mathbf4)$; upper and lower indices are exchanged as for the unitary groups $SU(2)_L$, $SU(2)_R$, $SU(4)_R$ of the compact real form. In Lorentzian signature, by contrast, $(\psi_\alpha)^\dagger$ is a dotted spinor. In ten-dimensional notation the conjugation maps $\zeta_A\in\mathbf{16}$ to an element of $\mathbf{16}'$, with upper index, compatibly with the gamma matrices: $\overline{\Gamma_M\zeta}=\Gamma_M\overline\zeta$ for the real basis vectors $e_M$.
\end{remark}

% ANCHOR: after_label thm:c8-dag
\begin{remark}[Physics notation: radial quantisation]
This is the conjugation of radial quantisation, $P^\dagger=K$, $Q^\dagger=S$, $\bar Q^\dagger=\bar S$, $D^\dagger=D$ in physics notation, here with the signs $P_b^\dagger=-K_{\bar b}$, $(Q^i_\alpha)^\dagger=-S^\alpha{}_i$, $(\bar Q_{\dot\alpha i})^\dagger=\bar S^{\dot\alpha i}$, $\Delta^\dagger=\Delta$, and with the matrix adjoint on $\mathfrak{su}(4)_R$. By the dilatation part $\tfrac12(\xi,\eta)\Delta$ of $[\xi,\rho(x)\eta]$, both $[Q^i_\alpha,(Q^i_\alpha)^\dagger]$ and $[\bar Q_{\dot\alpha i},(\bar Q_{\dot\alpha i})^\dagger]$ (no sum) contain $-\tfrac12\Delta=D$, the physics dilatation; (iii) states that on a superconformal primary (annihilated by $\Sigma^*$, i.e. by $S$ and $\bar S$) of scaling dimension $d>0$ the norms of its $Q$-descendants are positive for the sign $+$. The real form $\mathfrak{psu}(2,2|4)$ has even part $\mathfrak{su}(2,2)\oplus\mathfrak{su}(4)_R$, with $\mathfrak{su}(2,2)\cong\mathfrak{so}(4,2)$ the Lorentzian conformal algebra, and maximal compact even subalgebra $\mathfrak{su}(2)_L\oplus\mathfrak{su}(2)_R\oplus\mathfrak u(1)\oplus\mathfrak{su}(4)_R$, the $\mathfrak u(1)$ generated by $i\Delta$.
\end{remark}

% ANCHOR: after_label def:c8-Theta
\begin{remark}[Physics notation: Osterwalder--Schrader reflection]
$\Theta_n$ is the Osterwalder--Schrader reflection in the hyperplane orthogonal to $n$; for $n=e_4$ it is Euclidean time reflection $x^4\mapsto-x^4$, and $\dagger_n$ corresponds to the Hermitian conjugation of the Wick-rotated Lorentzian algebra. On the scalars it uses the reality of the $\mathbf6$, $\overline{A_{ij}}=A^{ij}$, and on the gauginos it exchanges $\lambda$ and $\bar\lambda$. For $n=e_4$, where $\sigma^4=i$,
\[
(Q^i_\alpha)^{\dagger_{e_4}}=-\bar Q_{\dot\alpha i}\big|_{\dot\alpha=\alpha},\qquad (\bar Q_{\dot\alpha i})^{\dagger_{e_4}}=-Q^i_\alpha\big|_{\alpha=\dot\alpha},
\]
the Euclidean form of the Lorentzian relation $(Q_\alpha)^\dagger=\bar Q_{\dot\alpha}$ between undotted and dotted supercharges. The reflection inner product $\langle a,b\rangle_n$ of Definition \ref{def:c8-inner} is the Osterwalder--Schrader inner product.
\end{remark}

% ANCHOR: end_of_section 8.8
\begin{remark}[Physics notation: BPZ conjugation]
$C$ is the conformal map from the half space to the unit ball that carries planar (Osterwalder--Schrader) quantisation to radial quantisation, and $\Theta_I$, the inversion $x\mapsto x/x^2$ with the conformal factors, is the BPZ conjugation of radial quantisation. Proposition \ref{prop:c8-primary} is the BPZ conjugate of a scalar primary of scaling dimension $w_O+\kappa$, with $\kappa$ its one-loop anomalous dimension ($D_1(O)=2\kappa O$),
\[
O(x)^\dagger=|x|^{-2(w_O+\kappa)}\,\overline O\big(x/|x|^2\big),
\]
the anomalous part $K_1(O)=-2(b\cdot x)\kappa O$ of the special conformal transformation being that of $b^m\big(x^2\partial_m-2x_m\,x\cdot\partial-2(w_O+\kappa)x_m\big)$.
\end{remark}

% ANCHOR: after_label prop:c8-2pt
\begin{remark}[Physics notation: two-point functions]
(i) is the decoupling of BRST-exact operators, (ii) the absence of ultraviolet and infrared problems at separated points, and (iii) the Callan--Symanzik equation of the two-point function, with $\ell$ the renormalisation length. For a single scalar operator with $D_1(O)=2\gamma O$ ($\gamma=\frac{3g^2N}{4\pi^2}$ for the Konishi operator, $0$ for the $\mathbf{20'}$) it reads
\[
\langle O,O\rangle(x,0)=|x|^{-2w}\Big[C+(\text{one-loop finite part})-\gamma\,C\log\frac{|x|^2}{\ell^2}\Big],
\]
the one-loop expansion of $C\,(|x|^2/\ell^2)^{-\gamma}|x|^{-2w}$, the power law of an operator of dimension $w+\gamma$.
\end{remark}

% ANCHOR: end_of_section 8.9
\begin{remark}[Physics notation: BPZ inner product and $Q^\dagger=S$]
The radial inner product is the BPZ inner product of radial quantisation, the two-point function of $b$ at the origin with the inversion conjugate of $a$, i.e. the overlap of the in-state of $b$ with the out-state of $a$. Its Gram matrix $\mathrm{Gram}_{\rm tree}+\hbar\,\mathrm{Gram}_1$ is the matrix of two-point coefficients of the operator basis, and $Y^{\rm adj}$ is the Hermitian adjoint for the tree metric. Corollary \ref{cor:c8-gram} is the one-loop form of $S=Q^\dagger$, $\bar S=\bar Q^\dagger$ and $K=P^\dagger$, with the signs of Theorem \ref{thm:c8-dag}: the one-loop special supercharges are the adjoints of the one-loop supercharges for the one-loop-corrected metric, the commutator with $\mathrm{Gram}_{\rm tree}^{-1}\mathrm{Gram}_1$ being the change of the metric. The relation for $D_1$ is the self-adjointness of the one-loop dilatation operator, hence the reality of the one-loop anomalous dimensions.
\end{remark}

% ANCHOR: end_of_section 8.10
\begin{remark}[Physics notation: higher loops]
Order $\hbar^k$ is $k$ loops, order $g^{2k}$ in the normalisation of Remark \ref{rem:c6-konishi}. At two loops the counterterm vertices of order $\hbar$ inserted into one-loop graphs subtract the one-loop subdivergences, and $L_2$ contains the two-loop corrections to all superconformal generators, $\tfrac12$ of its component along $\Delta$ being the two-loop dilatation operator.
\end{remark}
